\section{输出文件和后处理}

02 版程序输出 Tecplot ASCII 数据。输出函数是 \texttt{write\_tecplot()}。
文件开头的变量名为：
\begin{lstlisting}
VARIABLES = "x", "rho", "u", "p", "E",
            "rho_conserved", "rhou", "rhoE", "nu"
\end{lstlisting}

其中 \texttt{rho}、\texttt{u}、\texttt{p}、\texttt{E} 是原始量，
\texttt{rho\_conserved}、\texttt{rhou}、\texttt{rhoE} 是守恒量，
\texttt{nu} 是人工粘性滤波系数。

\subsection{数值解和精确解文件}

程序最终会输出一个 numerical 文件，并自动生成对应的 exact 文件。例如：
\[
\texttt{numerical\_02.dat}
\quad\rightarrow\quad
\texttt{numerical\_02\_exact.dat}.
\]

如果设置了 snapshot interval，程序也会在中间时刻输出快照，并且为每个快照调用一次精确解程序。

\subsection{与 Python 后处理的关系}

已有 Python 后处理脚本主要读取 \texttt{x}、\texttt{rho}、\texttt{u}、\texttt{p}。
02 版虽然多输出了几列，但这几个基本变量名仍然保留，所以按变量名读取的脚本仍然可以使用。

\subsection{人工粘性系数的观察}

新增的 \texttt{nu} 列可以用来观察人工粘性主要加在哪里。
一般来说，在光滑区域 \(\nu_i\) 应该比较小；
在激波或接触间断附近，二阶差分变大，\(\nu_i\) 也会变大。
这可以用来比较 \(\rho\)、\(u\)、\(p\) 三种开关变量对人工粘性分布的影响。
